\section{Introduction}
\label{sec:introduction}

\setlength{\columnsep}{6pt}
\begin{wrapfigure}{r}{0.44\textwidth}
\vspace{-4.2em}
\centering
\includegraphics[width=\linewidth,trim=0 72pt 0 0,clip]{teaser_revised.pdf}
\vspace{-0.45cm}
\caption{\textbf{Motivation.} Judgmental forecasts combine observations with
contextual knowledge about the process that generated them. Frontier models,
asked to forecast this question, answer with $P(H) = 0.68$ (see
App.~\ref{app:carnival-coin-probe}).}
\vspace{-1.0cm}
\label{fig:world-knowledge-evidence}
\end{wrapfigure}

Judgmental forecasting is a demanding form of reasoning central to decisions
about political, social, and economic futures. Unlike settings with abundant
historical data for statistical extrapolation, judgmental forecasting requires
integrating sparse, noisy evidence with contextual knowledge to form calibrated
beliefs about future outcomes
\citep{tetlock2016superforecasting,mellers2014psychological,lawrence2006judgmental}.
This makes judgmental forecasting fundamentally inductive: forecasters must
infer from limited observations the relationships that govern how evidence
impacts outcomes, and use those relationships to anticipate how events may
unfold
\citep{tenenbaum2006theory,griffiths2010probabilistic,tenenbaum2011mind,lake2017building}.
Figure~\ref{fig:world-knowledge-evidence} illustrates this identification problem
in a low-stakes setting. A researcher observes a carnival coin-flipping machine
and, after seeing two heads, must estimate the probability that the next flip
will be heads. A judgmental forecast may incorporate contextual knowledge (e.g.,
whether the game is rigged) to infer the
process generating the observed outcomes.

Language Models (LMs) now approach crowd accuracy in forecasting tasks
\citep{halawi2024approaching,schoenegger2024wisdom}, raising a deeper question
about what underlies their predictive accuracy. Prior work has shown that in some
settings, models recover latent structure that supports their prediction
\citep{li2023emergent}, while in others, they predict accurately despite learning
relationships that fail to generalize beyond their observed setting
\citep{vafa2025foundation}. This distinction is particularly consequential for
judgmental forecasting, where decision-makers often care not only about what will
happen, but how the forecast should change under alternative evidence or events.
Such counterfactual questions require models to capture predictive relationships
that generalize across changes in context, rather than merely reproduce accurate
probabilities in the observed setting.

\noindent
\textbf{Present work.}
We study whether LMs exhibit inductive reasoning in forecasting through
progressively stronger behavioral experiments spanning real-world prediction
markets and controlled synthetic environments. First, we test whether LMs react
selectively to evidence by studying whether they distinguish outcome-relevant
evidence from topically related but orthogonal information
(\hyperref[sec:exp1]{Exp.~1}). Second, we test whether LMs recover predictive
relationships from context by varying contextual information while holding
numerical evidence fixed in a controlled synthetic environment
(\hyperref[sec:exp2]{Exp.~2}). Finally, we test whether LMs transfer what they
learn to new settings (\hyperref[sec:exp3]{Exp.~3} and
\hyperref[sec:exp4]{Exp.~4}) by fine-tuning models on forecasting tasks and
evaluating transfer to unseen synthetic structures and held-out historical
prediction markets.

\noindent\textbf{Results.}
Across four experiments, LMs show behavioral signatures of inductive forecasting.
All nine systems update $1.9$--$7.5\times$ more to outcome-relevant than topical
evidence (\hyperref[sec:exp1]{Exp.~1}). In Coin City, our controlled task,
forecasts track context-selected relationships ($r=.52$--$.70$) and revise them
as direct evidence arrives; larger models recover episode-local conventions that
affect forecasts under intervention (\hyperref[sec:exp2]{Exp.~2}). Episode-matched
training improves joint domain-and-mechanism transfer beyond an answer-matched
shuffled control at all three tested Qwen3 sizes (8B MAE $4.46$ vs.~$4.71$;
\hyperref[sec:exp3]{Exp.~3}). On 318
held-out markets, proper-score training improves every tested base; several
trained checkpoints reach crowd performance (\hyperref[sec:exp4]{Exp.~4}).
Together, LMs select, revise, and transfer predictive relationships, and training
improves forecasts on unseen events.
